\begin{abstract}
Curvilinear structure segmentation is a challenging task: the targets are
often thin, low-contrast, and topologically fragile, while annotations are
scarce and test domains vary substantially.  Prompt-free adaptation of a
frozen SAM with lightweight adapters has made few-shot segmentation
practical, but existing decoders often produce multiple candidate masks and
rely on test-time gating or fusion to select among them.  This
input-dependent selection is difficult to calibrate under domain shift,
increases inference complexity, and can break structural connectivity.
These issues indicate that the bottleneck is not generating more
candidates, but how to deploy them.  We therefore propose Deterministic
Candidate Decoding (DCD), which replaces per-image selection with a fixed
deployment rule.  DCD trains four candidate heads with dense supervision
and complementary topology-aware objectives, but always deploys the same
fixed candidate head at inference.  This turns a multi-candidate decoder
into a deterministic single-pass predictor while retaining the benefits of
candidate specialization.  Trained with only 18 images from six Seen
datasets, DCD is evaluated on 12 datasets (8 Unseen and 4 Novel) spanning
cracks, retinal and coronary vessels, roads, wires, and tyres.  Against ten
representative CNN- and SAM-based baselines, DCD achieves
state-of-the-art macro-averaged performance on this 12-dataset few-shot
benchmark.  Averaged over four training seeds, it obtains 58.76 Dice,
45.17 IoU, 66.02 clDice, and 63.13 HD95, ranking first on all four macro
metrics among all compared methods.  These gains are obtained with only 18
training images and a single fixed candidate head at inference.  Although
the improvements are not perfectly uniform across every dataset, DCD
achieves the best result on every macro-averaged metric and remains
competitive on the datasets where it trades HD95 for overlap.
\end{abstract}
